\documentclass{article}

% Packages
\usepackage{geometry}
\usepackage{graphicx}
\usepackage{lipsum}
\usepackage[portuguese]{babel}

% Page setup
\geometry{a4paper, margin=2cm}

\begin{document}

\thispagestyle{empty} % Remove page number from first page

\begin{figure}[t]
    \includegraphics[width=3cm]{images/logo-puc-minas.png}
    \hspace{0.02\textwidth}
    \vline%
    \hspace{0.04\textwidth}
    \includegraphics[width=3cm]{images/logo-icei.jpeg}
\end{figure}

\hrulefill%
\vspace{\baselineskip}

\Large\noindent
\textbf{Pontifícia Universidade Católica de Minas Gerais} \\
\textbf{Instituto de Ciências Exatas e Informática} \\
\textbf{Departamento de Engenharia de Computação}

\begin{center}
    \vfill
    \Huge\textbf{Relatório: Trabalho Prático 2} \\
    \vspace{0.5cm}
    \Large\textbf{Registradores em VHDL} \\
    \vspace{1cm}
    \large \textbf{Professor(es)}: Antônio Hamilton Magalhães\\
    \vspace{0.5cm}
    \large \textbf{Aluno(s)}: Bruno Luiz Dias Alves de Castro \\
    \large \hspace{0.75cm} Rafael Ramos de Andrade \\
    \vfill
    \large Belo Horizonte \\ Campus Coração Eucarístico \\
    \vspace{\baselineskip}
    \large \today
\end{center}

\newpage
\thispagestyle{empty}
\tableofcontents
\newpage

\large

\section{Introdução}
Durante as aulas da disciplina de Sistemas Reconfiguráveis, fomos introduzidos à linguagem VHDL. VHDL (\textbf{V}HSIC \textbf{H}ardware \textbf{D}escription \textbf{L}anguage) é uma linguage de descrição de hardware. Com ela, podemos montar circuitos lógicos de maneira totalmente textual, o que da à linguagem um gigante vantagem ante à soluções visuais.

\subsection{TP1}
Como primeiro trabalho prático, são propostas montagens de dois circuitos utilizando a linguagem VHDL para a placa de sensolvimento Altera:

\begin{itemize}
    \item Um addr\_mux (Secção 2)
    \item Uma ULA (Unidade Lógica Aritimética) (Secção 3)
\end{itemize}

\subsubsection{Objetivos}

Entre os objetivos que temos com o desenvolvimento deste tarbalho prático podemos listar:

\begin{itemize}
    \item Aprender conceitos básicos da linguagem VHDL;
    \item Implementar utilizando programação concorrente os dois circuitos propostos.
    \item Compilar os circuitos e testar os resultados na placa de desenvolvimento Altera.
\end{itemize}

\section{ADDR\_MUX}
\lipsum[1]

\section{ULA (Unidade Lógica Aritimética)}
\lipsum[3]

\section{Conclusion}
\lipsum[4]

\end{document}